\Qhead{1.2}{Автоколебательная система с одной степенью свободы. Энергетические соотношения и методы расчёта}{Автоколебательная система с одной степенью свободы. Энергетические соотношения в автоколебательных системах. Методы расчета автоколебательных систем.}
\Subfile{1.2.1}{sections/01_oscillations/q02/s01}
\Subfile{1.2.2}{sections/01_oscillations/q02/s02}
\Subfile{1.2.3}{sections/01_oscillations/q02/s03}

\recap{Источник энергии, колебательная система и нелинейная обратная связь создают предельный цикл. Баланс определяет амплитуду, изменение баланса --- устойчивость. Выбор метода зависит от формы сигнала и требуемой задачи.}
\board{Петля генератора; $W_+(A)-W_-(A)$; цикл; уравнение Ван-дер-Поля и $\dot A$.}
\oral{Чем предельный цикл отличается от эллипса консервативного осциллятора?}
\bait{Можно оставить жёсткое возбуждение: неустойчивый цикл задаёт конечную затравку для перехода из покоя в колебания.}
\sources{Пятаков~\cite{pyatakov}, PDF-с.~53--61: устройство, томсоновский и релаксационный режимы; вопрос~2 конспекта коллеги --- основа последовательности объяснения. Энергетический расчёт и нормировка $2\mu$ выведены здесь.}
